\section{Conclusions and Future Work}
In this study, we presented \texttt{\named}, the largest and most diverse dataset of Indian legal cases, alongside \texttt{\namel}, a specialized language model fine-tuned for legal judgment prediction and explanation. Our findings demonstrate that domain-specific models, particularly those enhanced with legal data, significantly outperform generic large language models in both accuracy and quality of explanations provided. Notably, we achieved very good accuracy in the prediction task after including data from all court levels, underscoring the value of comprehensive datasets.

% Despite these advances, the study faced several challenges, including token and computational resource constraints, which limited our exploration of larger models' full potential. Additionally, the dataset's focus on English-language judgments excluded regional languages like Hindi and Bengali, highlighting a need for more inclusive data. The limited scale of expert evaluations due to resource constraints also suggests areas for future improvement.
% 
Future work will focus on expanding the dataset to include judgments in regional languages, better reflecting India's linguistic diversity. We plan to explore larger and more advanced models, potentially using more efficient quantization techniques and enhanced hardware resources, to better handle complex legal documents. Refining our fine-tuning methodologies by incorporating a broader range of legal documents, such as statutes and contracts, will further enrich the model's knowledge base.

% We will also focus on developing strategies to mitigate hallucinations in model outputs, a critical factor in legal applications where accuracy is paramount. Enhancing evaluation metrics to better capture the nuances of legal reasoning and the quality of explanations will provide a more comprehensive assessment of model performance. Expanding the scale of expert evaluations will offer deeper insights into the models' capabilities and areas for improvement.
By addressing these challenges and expanding the scope of our research, we aim to enhance the performance and reliability of AI models in the legal domain, contributing to more efficient and accurate legal decision-making processes.

% The code and datasets are available at: \ab{put url}.
